\documentclass[11pt]{article}

\usepackage[a4paper,margin=1in]{geometry}
\usepackage[T1]{fontenc}
\usepackage[utf8]{inputenc}
\usepackage{lmodern}
\usepackage{microtype}
\usepackage{setspace}
\usepackage{amsmath,amssymb}
\usepackage{booktabs,tabularx,array}
\usepackage{enumitem}
\usepackage[dvipsnames]{xcolor}
\usepackage{titlesec}
\usepackage{hyperref}

\hypersetup{
  colorlinks=true,
  linkcolor=black,
  urlcolor=MidnightBlue,
  citecolor=black,
  pdftitle={AKMOLA-DT-V1 Architecture Brief}
}

\definecolor{accent}{HTML}{1E4E79}
\definecolor{textgray}{HTML}{4A4A4A}

\titleformat{\section}{\Large\bfseries\color{accent}}{\thesection}{0.6em}{}
\titleformat{\subsection}{\normalsize\bfseries}{\thesubsection}{0.5em}{}

\setlength{\parindent}{0pt}
\setlength{\parskip}{0.55em}
\setstretch{1.08}
\renewcommand{\arraystretch}{1.2}

\newcommand{\vone}{\textsc{Mandatory in V1}}
\newcommand{\later}{\textsc{Deferred Beyond V1}}
\newcolumntype{Y}{>{\raggedright\arraybackslash}X}

\begin{document}

\begin{titlepage}
  \centering
  \vspace*{2.5cm}
  {\LARGE\bfseries AKMOLA-DT-V1\par}
  \vspace{0.35cm}
  {\Large A Regional Snowmelt-Aware Flood Intelligence Architecture\\for Akmola / Astana\par}
  \vspace{1.2cm}
  {\large Architecture Brief\par}
  \vspace{1.2cm}

  \begin{minipage}{0.88\textwidth}
    \small
    \color{textgray}
    This memorandum defines a buildable architecture for flood intelligence in weak-gradient steppe terrain. The design joins regional hydrologic context, snowmelt-first forcing, and nested local micro-relief simulation for priority exposed zones. It rejects generic AI-first framing and anchors credibility in historical event replay, boundary-aware storage logic, and impact-oriented outputs suitable for district decision workflows.
  \end{minipage}

  \vfill
  \begin{tabular}{rl}
    Project: & AKMOLA-DT-V1 \\
    Prepared for: & Regional flood intelligence development track \\
    Prepared by: & Student systems architecture team \\
    Date: & \today \\
  \end{tabular}
\end{titlepage}

\section{Executive Reset}

The earlier framing was technically ambitious but strategically exposed. It foregrounded algorithmic identity before hydrologic structure, and it treated a regionally coupled steppe system as if a local tile could stand on its own. That mismatch is easy to challenge in any serious review.

Four corrections were necessary. First, the project had to move away from generic AI flood language and return to its strongest scientific claim: flat-steppe micro-relief hydrology. Second, quantum and model-brand narratives had to be demoted; they are not the source of first-order reliability in an operationally credible V1. Third, regional boundary behavior had to be made explicit, because weak-gradient connectivity allows external water and antecedent storage to shape local outcomes. Fourth, forcing logic had to pivot from rainfall-first to snowmelt-first, reflecting the plausible seasonal flood regime in Akmola and Astana.

This reset is not cosmetic. It is a structural correction from a technically interesting concept to a defensible regional intelligence architecture.

\section{Corrected Central Thesis}

AKMOLA-DT-V1 is a regional snowmelt-aware flood intelligence architecture that combines Akmola-scale terrain and hydrologic context, boundary-aware depression storage and spill connectivity, and nested high-detail local simulation in priority flood-exposed zones; its credibility is established through historical event replay and impact-based validation rather than algorithmic novelty claims.

\section{Architectural Principles}

\begin{enumerate}[leftmargin=1.5em,itemsep=0.35em]
  \item \textbf{Regional context first.} Local predictions are constrained by upstream storage, boundary inflow, and network connectivity across the broader domain.
  \item \textbf{Snowmelt-first forcing.} Seasonal snow accumulation and melt release drive the primary event narrative; rainfall is a modifier, not the sole trigger.
  \item \textbf{Flat-steppe micro-relief sensitivity.} Weak gradients and shallow depressions control storage and threshold spill transitions.
  \item \textbf{Nested local windows.} High-detail simulation is concentrated where decisions are made: settlements, road corridors, and known hotspot basins.
  \item \textbf{Historical replay as a design constraint.} The system is designed around reconstructing known events, not added as a late evaluation convenience.
  \item \textbf{Impact-oriented outputs.} Deliverables must map to exposed roads, settlements, and district summaries, not only internal model states.
\end{enumerate}

\section{System Architecture}

Let $\Omega_R$ denote the regional domain, $G_R$ the regional hydro-context graph, $L_i$ local simulation windows, and $t$ event time. The architecture is organized into eight modules.

\subsection{Module 1 --- Regional Terrain and Hydro-Context Layer}
\textbf{What it does.} Builds the regional backbone: major channels, floodplains, wetlands, transport embankments, broad storage zones, and weak-gradient connectivity pathways.

\textbf{Why it exists.} Local flood behavior is not hydrologically closed in this setting.

\textbf{Why it matters in Akmola.} External inflow can arrive slowly through subtle gradients and still control local inundation timing.

\textbf{V1 status.} \vone.

\textbf{Formal role.}
\[
\mathcal{C}_R:\Omega_R \rightarrow G_R.
\]

\subsection{Module 2 --- DEM Conditioning / Hydro-Enforcement Layer}
\textbf{What it does.} Converts raw elevation $z_{\mathrm{raw}}$ into a hydrologically usable surface $z_{\mathrm{hydro}}$ by correcting artifacts while preserving true depressions and anthropogenic controls.

\textbf{Why it exists.} Flat-terrain predictions fail rapidly when false pits or digital barriers remain unresolved.

\textbf{Why it matters in Akmola.} Micro-relief is the governing signal; small vertical errors alter storage and spill topology.

\textbf{V1 status.} \vone.

\textbf{Formal role.}
\[
z_{\mathrm{raw}} \mapsto z_{\mathrm{hydro}}, \qquad
\text{subject to preservation of true storage basins.}
\]

\subsection{Module 3 --- Depression / Storage / Spill Graph Layer}
\textbf{What it does.} Constructs a threshold graph of fill-and-spill dynamics across local storage cells.

\textbf{Why it exists.} Flood propagation in flat terrain is governed by basin filling sequence and spill activation, not slope direction alone.

\textbf{Why it matters in Akmola.} Shallow depressions and saddle thresholds create non-linear connectivity changes during thaw periods.

\textbf{V1 status.} \vone.

\textbf{Formal role.}
\[
G_D=(V_D,E_D,\sigma,c,r),
\]
where $V_D$ are storage units, $E_D$ spill links, $\sigma$ spill elevations, $c$ capacities, and $r$ hydraulic resistance proxies.

\subsection{Module 4 --- Snowmelt and Hydrometeorological Forcing Layer}
\textbf{What it does.} Represents seasonal snow storage, thaw activation, rain-on-snow interaction, and baseline atmospheric forcing.

\textbf{Why it exists.} The forcing narrative must match regional event physics.

\textbf{Why it matters in Akmola.} A rain-only story is structurally weak for spring flood interpretation.

\textbf{V1 status.} \vone\ in simplified form.

\textbf{Formal role.}
\[
F_t=\{SWE_t,\;M_t,\;P_t,\;T_t,\;B_t\},
\]
with snow storage $SWE_t$, melt release $M_t$, rainfall $P_t$, thermal trigger $T_t$, and background/boundary forcing $B_t$.

\subsection{Module 5 --- Regional Connectivity / Boundary Inflow Layer}
\textbf{What it does.} Derives boundary conditions for each local window from regional state and forcing.

\textbf{Why it exists.} It prevents tile-isolation error and enforces cross-scale consistency.

\textbf{Why it matters in Akmola.} Weak gradients make antecedent storage and inflow histories decisive for local outcomes.

\textbf{V1 status.} \vone.

\textbf{Formal role.}
\[
B_{L_i,t}=\Gamma(G_R,F_t,X_{R,t}),
\]
where $B_{L_i,t}$ supplies inflow/outflow constraints for window $L_i$.

\subsection{Module 6 --- Local Simulation Window Layer}
\textbf{What it does.} Runs high-detail micro-relief flood simulation only in selected priority windows.

\textbf{Why it exists.} Full-region high-resolution dynamics are not required for credible decisions and are not feasible for V1.

\textbf{Why it matters in Akmola.} Local consequence mapping depends on preserving fine threshold geometry near settlements and roads.

\textbf{V1 status.} \vone\ for a limited set of windows.

\textbf{Formal role.}
\[
X_{L_i,t+\Delta t}=S_{L_i}\!\left(X_{L_i,t},F_{L_i,t},B_{L_i,t},G_{D,i}\right).
\]

\subsection{Module 7 --- Historical Event Replay / Validation Layer}
\textbf{What it does.} Reconstructs known events using archived forcing and compares predicted impacts against observed evidence.

\textbf{Why it exists.} Without replay, the system remains an unverified simulator.

\textbf{Why it matters in Akmola.} Stakeholder confidence depends on whether known flood patterns are reproduced in plausible order and extent.

\textbf{V1 status.} \vone.

\textbf{Formal role.}
\[
\mathcal{R}_e:\{F^{(e)}_{0:T_e},X^{(e)}_0\}\mapsto \{\hat{X}^{(e)}_t\}_{t=0}^{T_e},
\quad
\mathcal{M}_e=\mathcal{M}\!\left(\hat{X}^{(e)},Y^{(e)}_{\mathrm{hist}}\right).
\]

\subsection{Module 8 --- Output / Decision-Support Layer}
\textbf{What it does.} Produces hotspot maps, settlement and road exposure rankings, district summaries, and replay-ready event reports.

\textbf{Why it exists.} Authorities act on interpretable geospatial impact products, not latent model states.

\textbf{Why it matters in Akmola.} District-level planning requires traceable links between hazard states and infrastructure exposure.

\textbf{V1 status.} \vone.

\textbf{Formal role.}
\[
O_t=\Psi(\hat{X}_t, I, A),
\]
where $I$ denotes impact layers and $A$ administrative/infrastructure layers.

\section{Regional vs Local Logic}

This two-level design is the architectural center of gravity.

\textbf{Level 1: Regional backbone.} The regional layer tracks snow storage, melt progression, broad routing pathways, floodplain interaction, and zone-level risk transfer. Its purpose is context fidelity: where water is likely to originate, accumulate, and propagate across districts.

\textbf{Level 2: Nested local windows.} The local layer resolves micro-relief dynamics only in selected high-consequence areas. Its purpose is actionable detail: which roads are exposed, which settlement cells cross thresholds, and which basins transition from storage to spill.

Coupling is explicit. The regional layer provides $(B_{L_i,t},F_{L_i,t},R^{\mathrm{ctx}}_{L_i,t})$ to each window. Local windows return impact products $O^{\mathrm{imp}}_{L_i,t}$ to the regional decision layer. This avoids the two common failures of student systems: implausible whole-region fine-scale simulation and isolated tile models with boundary blind spots.

The result is disciplined rather than maximalist: regional realism where context dominates, local detail where consequences are judged.

\section{Snowmelt Pivot}

The forcing story must begin with seasonal storage and release:
\[
\text{Snow accumulation} \;\rightarrow\; \text{thaw activation} \;\rightarrow\; \text{regional routing/storage} \;\rightarrow\; \text{local overflow}.
\]

Rainfall remains relevant, especially during thaw transitions, but it is not the organizing driver for a serious Akmola spring flood architecture.

A concise V1 state representation is:
\[
X_t=\{h_t,q_t,s_t,SWE_t,\phi_t,b_t\},
\]
with surface depth $h_t$, lateral flux proxy $q_t$, memory state $s_t$, snow storage $SWE_t$, depression fill state $\phi_t$, and boundary state $b_t$.

Forcing can be expressed as:
\[
F_t=\{T_t,P_t,M_t,U_t\},
\]
with thermal forcing $T_t$, rainfall $P_t$, melt input $M_t$, and upstream/boundary input $U_t$.

V1 does not require full cryosphere physics. It requires a defensible seasonal storage-release mechanism with transparent assumptions and scenario bounds.

\section{Historical Validation Strategy}

Historical replay is mandatory because it tests the architecture against known consequences, not only internal consistency.

Evidence can include affected settlements, flooded road segments, district impact reports, documented hotspot zones, event timing windows, satellite-derived footprints where available, and archived snow-weather conditions.

\subsection*{Validation logic}
\begin{tabularx}{\textwidth}{@{}>{\bfseries}p{0.27\textwidth}Y@{}}
\toprule
Validation target & Practical interpretation \\
\midrule
Hotspot overlap & Did predicted high-risk zones intersect known flooded areas? \\
Settlement hit rate & Were affected villages or urban fringes correctly flagged? \\
Road exposure ranking & Were historically disrupted corridors ranked near the top? \\
District consistency & Did district-level ordering match observed impact patterns? \\
Onset timing & Did risk intensify in the right sequence and window? \\
False alarms & Did the system discriminate, or flood nearly everything? \\
\bottomrule
\end{tabularx}

\subsection*{Event set for V1}
A compact but credible event bank should include: one snowmelt-dominant spring flood, one mixed thaw-plus-rain event, one moderate localized event, and one weak or non-event case for false-positive control.

\subsection*{Metric posture}
Where full depth truth is unavailable, prioritize impact-grounded metrics: hotspot precision/recall, settlement and road hit rates, district rank correlation, onset timing error, and false-alarm behavior. If partial flood footprints exist, report IoU, F1, and CSI as supplementary map metrics.

\section{Minimum Viable Regional Product}

The minimum serious product is not a universal simulator. It is a regional spring-flood context engine with nested local consequence windows and replay-tested outputs.

\begin{tabularx}{\textwidth}{@{}>{\bfseries}p{0.34\textwidth}Y@{}}
\toprule
Required V1 component & Scope discipline \\
\midrule
Regional domain and hydro-context backbone & Moderate resolution is acceptable if connectivity is preserved \\
Conditioned DEM and hydro-enforced terrain & Must preserve true depressions and key barriers \\
Depression/spill graph & Core physical logic; non-negotiable \\
Simplified snowmelt-aware forcing & Seasonal storage-release with scenario controls \\
Nested local windows (3--5) & Focus on settlements, roads, and recurrent hotspots \\
Historical replay workflow & Event reconstruction and impact scoring \\
Decision outputs & Hotspot maps, exposure ranks, district summaries \\
\bottomrule
\end{tabularx}

Anything beyond this baseline---full-region high-detail simulation, real-time operational forecasting, full data assimilation, or advanced surrogate learning---belongs to later phases.

\section{Team Realism}

\textbf{Computer science track.} Geospatial ingestion, graph construction, window orchestration, simulation workflow engineering, output generation, and reproducible replay tooling.

\textbf{Applied mathematics track.} Storage-spill formalization, reduced dynamics, threshold calibration, stability checks, objective design for replay scoring, and sensitivity analysis.

\textbf{Policy and strategy track.} Stakeholder framing, validation case selection, impact taxonomy, communication discipline, and scope control for government-facing artifacts.

This partition is practical and non-romantic. The strongest student outcome comes from disciplined ownership boundaries with shared validation standards.

\section{Risk Register}

\begin{tabularx}{\textwidth}{@{}p{0.22\textwidth}Y Y@{}}
\toprule
\textbf{Risk} & \textbf{Why it is dangerous} & \textbf{Primary mitigation} \\
\midrule
DEM truth and conditioning errors
& Minor elevation artifacts can invert storage and spill behavior in flat terrain.
& Multi-stage conditioning audits, preservation checks for true depressions, and visual review of critical corridors. \\

False depressions and digital dams
& Artifacts can produce convincing but incorrect inundation patterns.
& Separate true storage from raster artifacts and anthropogenic barriers; enforce conservative correction rules. \\

Routing instability under weak gradients
& Small numerical inconsistencies can trigger large topology shifts.
& Use storage-first threshold logic and bounded reduced dynamics. \\

Incorrect local boundary assumptions
& Isolated windows can produce confident but physically invalid forecasts.
& Require explicit regional-to-local boundary coupling for every priority window. \\

Snowmelt forcing uncertainty
& Event timing and magnitude can shift if thaw assumptions are wrong.
& Use scenario ensembles and report uncertainty bounds in replay outputs. \\

Overclaiming novelty
& Inflated claims reduce institutional trust and technical credibility.
& Anchor novelty in architecture and validation workflow, not exotic computation labels. \\

Scope overbuild
& Team bandwidth gets consumed before a testable system exists.
& Freeze V1 around backbone, forcing pivot, nested windows, replay, and impact outputs. \\

AI-language overdependence
& Stakeholders may infer black-box risk without evidence of reliability.
& Lead with hydrology, replay performance, and interpretable exposure products. \\
\bottomrule
\end{tabularx}

\section{Final Recommendation}

\textbf{One-sentence project definition.} AKMOLA-DT-V1 should be built as a regional snowmelt-aware flood intelligence system that couples Akmola-scale hydrologic context with boundary-aware nested micro-relief simulation and validates performance through historical event replay and impact metrics.

\textbf{Best architecture path.} Build the regional backbone first, enforce terrain realism, implement depression-spill logic, add simplified snowmelt forcing, connect nested windows, and validate against historical events before considering advanced surrogates.

\textbf{Biggest mistake to avoid.} Do not build a visually impressive local AI simulator and present it as a regional flood system. Credibility will come from forcing realism, boundary discipline, and replay-tested impact outputs.

\end{document}
